\documentclass{beamer}
\usepackage{tikz,amsmath,hyperref,graphicx,stackrel,animate,media9,listings}
\usetikzlibrary{positioning,shadows,arrows,shapes,calc,dsp,chains}
\newcommand{\argmax}{\operatornamewithlimits{argmax}}
\newcommand{\argmin}{\operatornamewithlimits{argmin}}
\mode<presentation>{\usetheme{Frankfurt}}
\AtBeginSection[]
{
  \begin{frame}<beamer>
    \frametitle{Outline}
    \tableofcontents[currentsection,currentsubsection]
  \end{frame}
}
\title{Lecture 32: Series and parallel filter combinations}
\author{Mark Hasegawa-Johnson}
\date{ECE 401: Signal and Image Analysis}  
\begin{document}

% Title
\begin{frame}
  \maketitle
\end{frame}

% Title
\begin{frame}
  \tableofcontents
\end{frame}

%%%%%%%%%%%%%%%%%%%%%%%%%%%%%%%%%%%%%%%%%%%%
\section[Review]{Review: FIR and IIR Filters, and Transfer Functions}
\setcounter{subsection}{1}

\begin{frame}
  \frametitle{System Functions}
  A first-order autoregressive filter,
  \[
  y[n] = x[n]-bx[n-1]+ay[n-1],
  \]
  has the impulse response and transfer function
  \[
  h[n]=a^n u[n]-ba^{n-1}u[n-1] \leftrightarrow H(z)  = \frac{1-bz^{-1}}{1-az^{-1}},
  \]
  where $a$ is called the {\bf pole} of the filter, and $b$ is called
  its {\bf zero}.
\end{frame}

\begin{frame}
  \frametitle{Impulse Response of a Second-Order All-Pole Filter}

  A general second-order all-pole filter has the system function
  \begin{align*}
    H(z) = \frac{1}{(1-p_1z^{-1})(1-p_2z^{-1})}= \frac{1}{1-(p_1+p_2)z^{-1}+p_1p_2z^{-2}}
  \end{align*}
  If the filter is underdamped, then $p_2=p_1^*$, and we can write
  \[
  H(z) = \frac{1}{1-2e^{-\sigma_1}\cos(\omega_1)z^{-1}+ e^{-2\sigma_1}z^{-2}}
  \]
  and
  \[
  h[n] = \frac{e^{-\sigma_1(n+1)}\sin(\omega_1(n+1))}{e^{-\sigma_1}\sin(\omega_1)} u[n]
  \]
\end{frame}

%%%%%%%%%%%%%%%%%%%%%%%%%%%%%%%%%%%%%%%%%%%%
\section[Errors]{Effects of numerical errors on poles and zeros}

\begin{frame}[fragile]
  \frametitle{Transfer function of a general difference equation}

  Suppose you have a set of $N$ poles and $N$ zeros, including complex
  conjugate pairs, that you want to implement as a filter.  Your
  desired $H(z)$ is:
  \begin{align*}
    H(z) = \frac{\prod_{i=1}^N (1-z_iz^{-1})}{\prod_{i=1}^N (1-p_iz^{-1})}
  \end{align*}
  The first step is to multiply them out to find the coefficients:
  \begin{align*}
    H(z) = \frac{1-\sum_{k=1}^Nb_kz^{-k}}{1-\sum_{k=1}^N a_k z^{-k}}
  \end{align*}
  Then implement it as a difference equation:
  \begin{lstlisting}[language=Python]
    y[n] = x[n] - sum(b*x[n-1:n-N-1:-1])\\
    + sum(a*y[n-1:n-N-1:-1])
  \end{lstlisting}
\end{frame}

\begin{frame}
  \frametitle{What about errors?}

  Suppose now we've implemented the difference equation
  \begin{align*}
    y[n] = x[n]-\sum_{k=1}^Nb_kx[n-k]+\sum_{k=1}^Na_k y[n-k]
  \end{align*}
  But now, suppose there are some errors:
  \begin{itemize}
  \item Errors in $x[n]$ might be caused by noise, or by quantization
  \item Errors in $a_k$ might be caused by quantization, or by computation mistakes
  \end{itemize}
  How badly do those errors effect $y[n]$?
\end{frame}
  
\begin{frame}
  \frametitle{The effect of errors gets worse with higher-order filters}

  \begin{itemize}
  \item Suppose $x[n]$ has some errors.  If those errors are near a pole of
    $H(z)$, then $y[n]$ will have exponentially worse errors.
    \begin{itemize}
    \item The number of poles gets larger as $N$ gets larger.
    \end{itemize}
  \item Suppose $a_k$ has an error.  Its effect on the poles, $p_i$, can
    be estimated using something called the implicit function theorem.
    \begin{itemize}
    \item The effect of each error, and the probability of a large error,
      increases exponentially as $N$ increases.
    \end{itemize}
  \item Both of these effects are much worse for $a_k$ than for $b_k$, because
    poles increase the signal amplitude, whereas zeros decrease it.
  \end{itemize}
\end{frame}

\begin{frame}
  \frametitle{Reducing errors $\Leftrightarrow$ Decompose the filter into sections}

  \begin{itemize}
  \item Error increases exponentially with $N$
  \item We can reduce the dependence to linear by decomposing the filter
    into a number of small sub-filters:
    \begin{align*}
      H(z) = H_1(z)+H_2(z)+\cdots + H_K(z)
    \end{align*}
    Then
    \begin{align*}
      y[n] = \left(h_1[n]\ast x[n]\right)+\cdots+\left(h_K[n]\ast x[n]\right)
    \end{align*}
    So the total error in $y[n]$ is just the sum of the errors in the
    sub-filters, i.e., linear dependence.
  \item Let's talk about how to do that decomposition.
  \end{itemize}
\end{frame}
  


%%%%%%%%%%%%%%%%%%%%%%%%%%%%%%%%%%%%%%%%%%%%
\section[Series]{Series connection of second-order sections}

\begin{frame}
  \frametitle{Why use block diagrams?}

  A first-order difference equation looks like
  \[
  y[n] = b_0x[n]-b_1x[n-1] +ay[n-1]
  \]
  \begin{itemize}
  \item It's pretty easy to understand what computation is taking
    place in a first-order difference equation.
  \item As we get to higher-order systems, though, the equations for
    implementing them will be kind of complicated.
  \item In order to make the complicated equations very easy, we
    represent the equations using block diagrams.
  \end{itemize}
\end{frame}

\begin{frame}
  \frametitle{Elements of a block diagram}

  A block diagram has just three main element types:
  \begin{enumerate}
  \item {\bf Multiplier:} the following element means
    $y[n]=b_0x[n]$:
    \vspace*{5mm}
    
    \centerline{\begin{tikzpicture}
        \node[dspnodeopen,dsp/label=right] (m02) at (2,0) {$y[n]$};
        \node[dspmixer,dsp/label=above] (m01) at (1,0) {$b_0$} edge[dspflow] (m02);
        \node[dspnodeopen,dsp/label=left] (m00) at (0,0) {$x[n]$} edge[dspflow](m01);
    \end{tikzpicture}}
  \item {\bf Unit Delay:} the following element means $y[n]=x[n-1]$
    (i.e., $Y(z)=z^{-1}X(z)$):
    \vspace*{3mm}
    
    \centerline{\begin{tikzpicture}
        \node[dspnodeopen,dsp/label=right] (m02) at (2,0) {$y[n]$};
        \node[dspsquare] (m01) at (1,0) {$z^{-1}$}  edge[dspflow] (m02);
        \node[dspnodeopen,dsp/label=left] (m00) at (0,0) {$x[n]$} edge[dspflow](m01);
    \end{tikzpicture}}          
  \item {\bf Adder:} the following element means $z[n]=x[n]+y[n]$:
    \vspace*{3mm}
    
    \centerline{\begin{tikzpicture}
        \node[dspnodeopen,dsp/label=right] (m12) at (2,-1) {$z[n]$};
        \node[dspadder](m11) at (1,-1) {}  edge[dspflow] (m12);
        \node[dspnodeopen,dsp/label=left] (m10) at (0,-1) {$y[n]$} edge[dspflow](m11);
        \node[coordinate](m01) at (1,0) {}  edge[dspflow](m11);
        \node[dspnodeopen,dsp/label=left] (m00) at (0,0) {$x[n]$} edge[dspflow](m01);
    \end{tikzpicture}}
  \end{enumerate}
\end{frame}

\begin{frame}
  \frametitle{Example: Time Domain}

  Here's an example of a complete block diagram:
  \centerline{\begin{tikzpicture}
      \node[dspnodeopen,dsp/label=right] (y) at (6,2) {$y[n]$};
      \node[dspsquare] (ydelay) at (5,1) {$z^{-1}$};
      \node[dspadder] (adder) at (3,2) {};
      \node[dspnodefull](ysplit) at (5,2){} edge[dspflow](y) edge[dspconn](ydelay) edge[dspline](adder);
      \node[coordinate] (ycenter) at (3,0) {} edge[dspconn] (adder);
      \node[dspmixer,dsp/label=above] (ymix) at (4,0) {$a$} edge[dspline] (ycenter);
      \node[coordinate] (ycorner)  at (5,0) {} edge[dspconn] (ymix) edge[dspline](ydelay);
      \node[dspnodeopen,dsp/label=left] (x) at (2,2) {$x[n]$} edge[dspconn](adder);
  \end{tikzpicture}}

  This block diagram is equivalent to the following equation:
  \begin{displaymath}
    y[n] = x[n]+ay[n-1]
  \end{displaymath}
  Notice that we can read it, also, as
  \begin{displaymath}
    Y(z)  = X(z) + az^{-1}Y(z)~~~\Rightarrow~~~H(z)=\frac{1}{1-az^{-1}}
  \end{displaymath}
\end{frame}

\begin{frame}
  \frametitle{A Complete First-Order IIR Filter}

  Now consider how we can represent a complete first-order IIR filter, including both
  the pole and the zero.  Here it is in the $z$-domain:
  \begin{displaymath}
    Y(z)  = b_0X(z) -b_1z^{-1}X(z)+ a_1z^{-1}Y(z).
  \end{displaymath}
  When we implement it, we would write a line of python that does this:
  \begin{displaymath}
    y[n] = b_0x[n]-b_1x[n-1]+a_1y[n-1],
  \end{displaymath}
  which is exactly this block diagram:
  \vspace*{3mm}
  
  \centerline{\begin{tikzpicture}
      \node[dspnodeopen,dsp/label=right] (y) at (6,2) {$y[n]$};
      \node[dspsquare] (ydelay) at (5,1) {$z^{-1}$};
      \node[dspsquare] (xdelay) at (1,1) {$z^{-1}$};
      \node[dspadder] (adder) at (3,2) {};
      \node[dspmixer,dsp/label=below] (b0) at (2,2) {$b_0$} edge[dspconn] (adder);
      \node[dspnodefull](ysplit) at (5,2){} edge[dspflow](y) edge[dspconn](ydelay) edge[dspline](adder);
      \node[dspnodefull] (xsplit) at (1,2) {} edge[dspconn] (xdelay) edge[dspline](b0);
      \node[coordinate] (ycenter) at (3.1,0) {} edge[dspconn] (adder);
      \node[coordinate] (xcenter) at (2.9,0) {} edge[dspconn] (adder);
      \node[dspmixer,dsp/label=above] (ymix) at (4,0) {$a_1$} edge[dspline] (ycenter);
      \node[dspmixer,dsp/label=above] (xmix) at (2,0) {$-b_1$} edge[dspline] (xcenter);
      \node[coordinate] (ycorner)  at (5,0) {} edge[dspconn] (ymix) edge[dspline](ydelay);
      \node[coordinate] (xcorner)  at (1,0) {} edge[dspconn] (xmix) edge[dspline](xdelay);
      \node[dspnodeopen,dsp/label=left] (x) at (0,2) {$x[n]$} edge[dspline](xsplit);
  \end{tikzpicture}}
\end{frame}

\begin{frame}
  \frametitle{Series and Parallel Combinations}
  Now let's talk about how to combine systems.
  \begin{itemize}
  \item {\bf Series combination}: passing the signal through two
    systems {\bf in series} is like multiplying the system functions:
    \[
    H(z)=H_2(z)H_1(z)
    \]
  \item {\bf Parallel combination}: passing the signal through two
    systems in {\bf parallel}, then adding the outputs, is like adding
    the system functions:
    \[
    H(z) = H_1(z)+H_2(z)
    \]
  \end{itemize}
\end{frame}
  

\begin{frame}
  \frametitle{One Block for Each System}

  Suppose that one of the two systems, $H_1(z)$, looks like this:
  \vspace*{3mm}

  \centerline{\begin{tikzpicture}
      \node[dspnodeopen,dsp/label=right] (y) at (6,2) {$y[n]$};
      \node[dspsquare] (ydelay) at (5,1) {$z^{-1}$};
      \node[dspadder] (adder) at (3,2) {};
      \node[dspnodefull](ysplit) at (5,2){} edge[dspflow](y) edge[dspconn](ydelay) edge[dspline](adder);
      \node[coordinate] (ycenter) at (3,0) {} edge[dspconn] (adder);
      \node[dspmixer,dsp/label=above] (ymix) at (4,0) {$p_1$} edge[dspline] (ycenter);
      \node[coordinate] (ycorner)  at (5,0) {} edge[dspconn] (ymix) edge[dspline](ydelay);
      \node[dspnodeopen,dsp/label=left] (x) at (2,2) {$x[n]$} edge[dspconn](adder);
  \end{tikzpicture}}
  and has the system function
  \begin{displaymath}
    H_1(z) = \frac{1}{1-p_1z^{-1}}
  \end{displaymath}
  Let's represent the whole system using a single box:
  \vspace*{3mm}

  \centerline{\begin{tikzpicture}
      \node[dspnodeopen,dsp/label=right] (m02) at (2,0) {$y[n]$};
      \node[dspsquare] (m01) at (1,0) {$H_1(z)$}  edge[dspflow] (m02);
      \node[dspnodeopen,dsp/label=left] (m00) at (0,0) {$x[n]$} edge[dspflow](m01);
  \end{tikzpicture}}          
\end{frame}

\begin{frame}
  \frametitle{Series Combination}

  The series combination, then, looks like this:
  \vspace*{3mm}
  
  \centerline{\begin{tikzpicture}
      \node[dspnodeopen,dsp/label=right] (y2) at (2,0) {$y[n]$};
      \node[dspsquare] (h2) at (1,0) {$H_2(z)$}  edge[dspconn] (y2);
      \node[dspnodefull,dsp/label=above] (y1) at (0,0) {$v[n]$} edge[dspconn] (h2);
      \node[dspsquare] (h1) at (-1,0) {$H_1(z)$}  edge[dspline] (y1);
      \node[dspnodeopen,dsp/label=left] (x) at (-2,0) {$x[n]$} edge[dspconn] (h1);
  \end{tikzpicture}}
  This means that
  \[
  Y(z) = H_2(z)V(z) = H_2(z)H_1(z)X(z)
  \]
  and therefore
  \[
  H(z) = \frac{Y(z)}{X(z)} = H_1(z)H_2(z)
  \]
\end{frame}

\begin{frame}
  \frametitle{Series Combination}

  The series combination, then, looks like this:
  \vspace*{3mm}

  \centerline{\begin{tikzpicture}
      \node[dspnodeopen,dsp/label=right] (y2) at (2,0) {$y_2[n]$};
      \node[dspsquare] (h2) at (1,0) {$H_2(z)$}  edge[dspconn] (y2);
      \node[coordinate,dsp/label=above] (y1) at (0,0) {$y_1[n]$} edge[dspconn] (h2);
      \node[dspsquare] (h1) at (-1,0) {$H_1(z)$}  edge[dspline] (y1);
      \node[dspnodeopen,dsp/label=left] (x) at (-2,0) {$x[n]$} edge[dspconn] (h1);
  \end{tikzpicture}}
  Suppose that we know each of the systems separately:
  \[
  H_1(z)=\frac{1}{1-p_1z^{-1}},~~~~~
  H_2(z)=\frac{1}{1-p_2z^{-1}}
  \]
  Then, to get $H(z)$, we just  have to multiply:
  \[
  H(z) = \frac{1}{(1-p_1z^{-1})(1-p_2z^{-1})} =
  \frac{1}{1-(p_1+p_2)z^{-1}+p_1p_2z^{-2}}
  \]
\end{frame}


%%%%%%%%%%%%%%%%%%%%%%%%%%%%%%%%%%%%%%%%%%%%
\section[Parallel]{Parallel connection of second-order sections}

\begin{frame}
  \frametitle{Parallel Combination}

  Parallel combination of two systems looks like this:
  \vspace*{3mm}

  \centerline{\begin{tikzpicture}
      \node[dspnodeopen,dsp/label=right] (y) at (2,0) {$y[n]$};
      \node[dspadder] (adder) at (1,0) {}  edge[dspflow] (y);
      \node[coordinate] (y1) at (1,1) {}  edge[dspline] (adder);
      \node[coordinate] (y2) at (1,-1) {}  edge[dspline] (adder);
      \node[dspsquare] (h1) at (0,1) {$H_1(z)$}  edge[dspline] (y1);
      \node[dspsquare] (h2) at (0,-1) {$H_2(z)$}  edge[dspline] (y2);
      \node[coordinate] (x1) at (-1,1) {}  edge[dspconn] (h1);
      \node[coordinate] (x2) at (-1,-1) {}  edge[dspconn] (h2);
      \node[dspnodefull] (xsplit) at (-1,0) {} edge[dspline](x1) edge[dspline](x2);
      \node[dspnodeopen,dsp/label=left] (x) at (-2,0) {$x[n]$} edge[dspline] (xsplit);
  \end{tikzpicture}}
  This means that
  \[
  Y(z) = H_1(z)X(z)+H_2(z)X(z)
  \]
  and therefore
  \[
  H(z) = \frac{Y(z)}{X(z)} = H_1(z)  + H_2(z)
  \]
\end{frame}

\begin{frame}
  \frametitle{Parallel Combination}

  Parallel combination of two systems looks like this:
  \vspace*{3mm}

  \centerline{\begin{tikzpicture}
      \node[dspnodeopen,dsp/label=right] (y) at (2,0) {$y[n]$};
      \node[dspadder] (adder) at (1,0) {}  edge[dspflow] (y);
      \node[coordinate] (y1) at (1,1) {}  edge[dspline] (adder);
      \node[coordinate] (y2) at (1,-1) {}  edge[dspline] (adder);
      \node[dspsquare] (h1) at (0,1) {$H_1(z)$}  edge[dspline] (y1);
      \node[dspsquare] (h2) at (0,-1) {$H_2(z)$}  edge[dspline] (y2);
      \node[coordinate] (x1) at (-1,1) {}  edge[dspconn] (h1);
      \node[coordinate] (x2) at (-1,-1) {}  edge[dspconn] (h2);
      \node[dspnodefull] (xsplit) at (-1,0) {} edge[dspline](x1) edge[dspline](x2);
      \node[dspnodeopen,dsp/label=left] (x) at (-2,0) {$x[n]$} edge[dspline] (xsplit);
  \end{tikzpicture}}
  Suppose that we know each of the systems separately:
  \[
  H_1(z)=\frac{1}{1-p_1z^{-1}},~~~~~
  H_2(z)=\frac{1}{1-p_2z^{-1}}
  \]
  Then, to get $H(z)$, we just  have to add:
  \[
  H(z) = \frac{1}{1-p_1z^{-1}}+\frac{1}{1-p_2z^{-1}}
  = \frac{2-(p_1+p_2)z^{-1}}{1-(p_1+p_2)z^{-1}+p_1p_2z^{-2}}
  \]
\end{frame}

\begin{frame}
  \frametitle{Second-order sections}

  \begin{itemize}
  \item Until now, we've pretended that you want to decompose $H(z)$
    into first-order sections.
  \item That's never true, because if $p_1$ is complex, then the first-order filter has
    complex-valued coefficients:
    \begin{align*}
      y[n] = x[n]+p_1 y[n-1]
    \end{align*}
  \item Real implementations like \texttt{scipy.signal.sosfilt}
    decompose an $N^{\text{th}}$-order filter into $\frac{N}{2}$
    second-order sections.  Each section has poles $p_2=p_1^*$, so its
    coefficients are real-valued:
    \begin{align*}
      y[n] = x[n]-by[n-1]+cy[n-2]
    \end{align*}
    where $b=2\Re(p_1)$ is real, and $c=|p_1|^2$ is also real.
  \end{itemize}
\end{frame}


%\begin{frame}
%  \frametitle{Parallel Combination}
%
%  Parallel combination of two systems looks like this:
%  \vspace*{3mm}
%
%  \centerline{\begin{tikzpicture}
%      \node[dspnodeopen,dsp/label=above] (y) at (2,0) {$y[n]$};
%      \node[dspadder] (adder) at (1,0) {}  edge[dspflow] (y);
%      \node[coordinate] (y1) at (1,1) {}  edge[dspline] (adder);
%      \node[coordinate] (y2) at (1,-1) {}  edge[dspline] (adder);
%      \node[dspsquare] (h1) at (0,1) {$H_1(z)$}  edge[dspline] (y1);
%      \node[dspsquare] (h2) at (0,-1) {$H_2(z)$}  edge[dspline] (y2);
%      \node[coordinate] (x1) at (-1,1) {}  edge[dspconn] (h1);
%      \node[coordinate] (x2) at (-1,-1) {}  edge[dspconn] (h2);
%      \node[dspnodefull] (xsplit) at (-1,0) {} edge[dspline](x1) edge[dspline](x2);
%      \node[dspnodeopen,dsp/label=above] (x) at (-2,0) {$x[n]$} edge[dspline](xsplit);
%  \end{tikzpicture}}
%  \begin{align*}
%  H(z) &= \frac{1}{1-p_1z^{-1}}+\frac{1}{1-p_2z^{-1}}\\
%  &= \frac{1-p_2z^{-1}}{(1-p_1z^{-1})(1-p_2z^{-1})}+\frac{1-p_1z^{-1}}{(1-p_1z^{-1})(1-p_2z^{-1})}\\
%  &= \frac{2-(p_1+p_2)z^{-1}}{1-(p_1+p_2)z^{-1}+p_1p_2z^{-2}}
%  \end{align*}
%\end{frame}


%%%%%%%%%%%%%%%%%%%%%%%%%%%%%%%%%%%%%%%%%%%%
\section[PFE]{Partial fraction expansion}

\begin{frame}
  \frametitle{Parallel decomposition}

  Suppose you know $H(z)$, and you want to find $H_1(z),\ldots,H_N(z)$ so that
  \begin{align*}
    H(z) &= H_1(z)+\cdots + H_N(z)
  \end{align*}
  How can we find $H_1(z)$ through $H_N(z)$?
\end{frame}
      
\begin{frame}
  \frametitle{Step \#1: Numerator Terms}

  Suppose you already know that
  \begin{displaymath}
    H(z) = \frac{1-\sum_{k=1}^N b_kz^{-k}}{A(z)} = G(z)\left(1-\sum_{k=1}^N b_kz^{-k}\right)
  \end{displaymath}
  and suppose you already know $g[n]$, the inverse transform of $G(z)=1/A(z)$.  Then
  $h[n]$ is just
  \begin{displaymath}
    h[n] = g[n] -\sum_{k=1}^N b_k g[n-k]
  \end{displaymath}
\end{frame}

\begin{frame}
  \frametitle{Step \#1: Numerator Terms}

  For example, here is the inverse transform of $H(z)=\frac{1+0.5z^{-1}}{1-0.85z^{-1}}$:
  \centerline{\includegraphics[width=4.5in]{exp/numsum.png}}
\end{frame}

\begin{frame}
  \frametitle{Step \#2: Denominator Terms}

  Now we need to figure out how to decompose
  \begin{displaymath}
    G(z) = \frac{1}{A(z)}
  \end{displaymath}
  We will solve this using  a method called {\bf partial fraction expansion.}
\end{frame}

\begin{frame}
  \frametitle{Step \#2: Partial Fraction Expansion}
  Partial fraction expansion works like this:
  \begin{enumerate}
  \item Factor $A(z)$:
    \begin{displaymath}
      G(z) = \frac{1}{\prod_{\ell=1}^N \left(1-p_\ell z^{-1}\right)}
    \end{displaymath}
  \item Assume that $G(z)$ is the result of a parallel system
    combination:
    \begin{displaymath}
      G(z) = \frac{C_1}{1-p_1z^{-1}} + \frac{C_2}{1-p_2z^{-1}} + \cdots
    \end{displaymath}
  \item Find the constants, $C_\ell$, that make the equation true.
    Such constants always exist, as long as none of the roots are
    repeated ($p_k\ne p_\ell$ for $k\ne\ell$).
  \end{enumerate}
\end{frame}

\begin{frame}
  \frametitle{Partial Fraction Expansion: Example}
  Step \# 1:  Factor it:
  \begin{displaymath}
    \frac{1}{1-1.2z^{-1}+0.72z^{-2}}=
    \frac{1}{\left(1-(0.6+j0.6)z^{-1}\right)\left(1-(0.6-j0.6)z^{-1}\right)}
  \end{displaymath}
  Step \#2: Express it as a  sum:
  \begin{displaymath}
    \frac{1}{1-1.2z^{-1}+0.72z^{-2}}=
    \frac{C_1}{1-(0.6+j0.6)z^{-1}}+\frac{C_2}{1-(0.6-j0.6)z^{-1}}
  \end{displaymath}
  Step \#3: Find the constants.
\end{frame}

\begin{frame}
  \frametitle{Partial Fraction Expansion: Example}
  The system function is:
  \begin{align*}
    G(z) &= \frac{1}{1-1.2z^{-1}+0.72z^{-2}}\\
    &=\frac{C_1}{1-(0.6+j0.6)z^{-1}}+\frac{C_2}{1-(0.6-j0.6)z^{-1}}
  \end{align*}
  and therefore the impulse response is:
  \begin{align*}
    g[n] &= C_1(0.6+0.6j)^nu[n]+C_2(0.6-j0.6)^nu[n]
  \end{align*}
  After finding the constants and simplifying, we get
  \begin{align*}
    g[n]&= \sqrt{2}(0.6\sqrt{2})^n \cos\left(\frac{\pi}{4}(n-1)\right)u[n]
  \end{align*}
\end{frame}

\begin{frame}
  \centerline{\includegraphics[width=4.5in]{exp/densum.png}}
\end{frame}

\begin{frame}
  \frametitle{How to find the inverse Z transform}

  Any IIR filter $H(z)$ can be written as\ldots
  \begin{itemize}
  \item a {\bf partial fraction expansion} into a sum of {\bf
    denominator} terms, each with this form:
    \begin{displaymath}
      G_\ell(z)=\frac{1}{1-az^{-1}}~~~\leftrightarrow~~~g_\ell[n]= a^nu[n],
    \end{displaymath}
  \item each possibly multiplied by a {\bf numerator} term, like this one:
    \begin{displaymath}
      D_k(z)=b_kz^{-k}~~~\leftrightarrow~~~d_k[n]=b_k\delta[n-k].
    \end{displaymath}
  \end{itemize}
\end{frame}

\begin{frame}
  \frametitle{How to implement second-order sections}

  Any IIR filter $H(z)$ can be written as\ldots
  \begin{itemize}
  \item a {\bf partial fraction expansion} into a sum of {\bf
    denominator} terms, each with this form:
    \begin{displaymath}
      G_\ell(z)=\frac{1}{1-az^{-1}}~~~\leftrightarrow~~~g_\ell[n]= a^nu[n],
    \end{displaymath}
    \ldots which you then combine in complex-conjugate pairs, to
    create $\frac{N}{2}$ second-order sections with real coefficients.
  \item each possibly multiplied by a {\bf numerator} term, like this one:
    \begin{displaymath}
      D_k(z)=b_kz^{-k}~~~\leftrightarrow~~~d_k[n]=b_k\delta[n-k].
    \end{displaymath}
  \end{itemize}
\end{frame}

%%%%%%%%%%%%%%%%%%%%%%%%%%%%%%%%%%%%%%%%%%%%
\section[Summary]{Summary}
\setcounter{subsection}{1}

\begin{frame}
  \frametitle{Summary: Block Diagrams}
  \begin{itemize}
  \item A {\bf block diagram} shows the delays, additions, and
    multiplications necessary to compute output from input.
  \item {\bf Series combination}: passing the signal through two
    systems {\bf in series} is like multiplying the system functions:
    \[
    H(z)=H_2(z)H_1(z)
    \]
  \item {\bf Parallel combination}: passing the signal through two
    systems in {\bf parallel}, then adding the outputs, is like adding
    the system functions:
    \[
    H(z) = H_1(z)+H_2(z)
    \]
  \end{itemize}
\end{frame}
  
\begin{frame}
  \frametitle{Summary: Inverse Z Transform}

  Any IIR filter $H(z)$ can be written as\ldots
  \begin{itemize}
  \item a {\bf partial fraction expansion} into a sum of {\bf
    denominator} terms, each with this form:
    \begin{displaymath}
      G_\ell(z)=\frac{1}{1-az^{-1}}~~~\leftrightarrow~~~g_\ell[n]= a^nu[n],
    \end{displaymath}
    \ldots which you can then combine in complex-conjugate pairs if you want
    difference equations with real-valued coefficients.
  \item each possibly multiplied by a {\bf numerator} term, like this one:
    \begin{displaymath}
      D_k(z)=b_kz^{-k}~~~\leftrightarrow~~~d_k[n]=b_k\delta[n-k].
    \end{displaymath}
  \end{itemize}
\end{frame}


\end{document}
